\documentclass[12pt, oneside, letterpaper]{book}
%%
%
%%
%%% DOCUMENT START
%%%General
\usepackage{verbatim} % For comments and other
\usepackage{framed}
\usepackage{multicol}
%%%
%%
%%% Subject Specific
%% Math
\usepackage{mathtools}
\usepackage{amsmath}  % For math
\allowdisplaybreaks % Allows page breaks inside long math blocks
\usepackage{amssymb}  % For more math
\usepackage{siunitx}  % standardization of si units
\DeclareSIUnit\Calorie{Cal}
\DeclareSIUnit\calorie{cal}
\DeclareSIUnit\molc{molecules}
\DeclareSIUnit\atoms{atoms}
\DeclareSIUnit\Molar{M}
\usepackage{physunits}
\usepackage{pgf}

\usepackage{unicode-math} % to ensure compatibility to correctly display math symbols
%%
%
%% Chemistry
\usepackage[version=4]{mhchem} % Package to help with formatting chemical formulae
%%
\usepackage{modiagram} % Allows drawing molecular orbital diagrams
%
\setmodiagram{style = fancy, distance = 7cm, AO-width = 15pt, labels, names} % Setting options for the molecular orbital diagrams
%%
\usepackage{chemmacros}
%%
%
%% Physics
\usepackage{physics2}  % for simplified dv, and v
%%
%
%%
%%% Graphics, tables, figures
%% Graphics
\usepackage{graphicx} % Required for including pictures
\usepackage[figurename=Figure]{caption}
\usepackage{subfig}   % For subfigures
\usepackage{xcolor,cancel}
\renewcommand{\CancelColor}{\color{red}}
\usepackage{wrapfig}
%%
%
%% Making Figures
\usepackage{tikz,bm} % Useful for drawing plots
\usepackage{tikz-3dplot}
\usepackage{circuitikz}
%%
%
%% Tables
\usepackage{float} % For tables and other floats
\usepackage[table]{xcolor} % Setting parameters for how tables will be formatted
\usepackage{tabularx}
\usepackage{array}
\usepackage{booktabs} % For making nicer tables
%%
%
%%
%%% General Formatting
\usepackage{ragged2e} %text alignment commands
\usepackage{fullpage} % Set margins and place page numbers at bottom center
\usepackage{paralist} % paragraph spacing
\usepackage{listings} % For source code
\usepackage{fancyhdr} % Headers and Footers
\usepackage{enumitem} % Lists
\usepackage{url} % url formatting
\usepackage{titlesec}
%%%
%%
%%% Fonts and characters
\usepackage{fontspec}
%%%
%%
%%% Setting full document settings
\setlength{\arrayrulewidth}{0.5mm}
\setlength{\tabcolsep}{20pt}
\usepackage[paperheight=279.4mm, paperwidth=215.9mm,
nomarginpar,% We don't want any margin paragraphs
textwidth = 177.7mm,
tmargin = 25.4mm,
bmargin = 25.4mm,
headheight = 10mm]% Set \headheight to 10mm%]% Set \headheight to 10mm
{geometry}
%% Setting up titles for sections and subsections. Setting up paragraph formatting
\newcommand{\styleTitle}[1] {\begin{flushleft}
       \textbf{\Large{#1}}
    \end{flushleft}
\vspace{1mm}
	\hrule}

\titleformat{\section}
{\center\large\bfseries}{}{1em}{}

\titleformat{\subsection}
{\RaggedRight\large\bfseries}{}{1em}{}

\titleformat{\subsubsection}
{\RaggedRight\normalsize\bfseries}{}{1em}{}

\titleformat{\paragraph}
{\normalfont\normalsize\mdseries}{}{0em}{}
%%
%
%% Space between paragraphs
\usepackage[skip=12pt plus1pt, indent=0pt]{parskip}
%% % Moved to another file for ease of updating
\begin{document}
%%
%
%%
%%% Document Setup Settings
%%
\pagestyle{fancy}
\renewcommand{\headrulewidth}{0pt}
\fancyhead{}
\fancyhead[R]{\thepage}
\fancyfoot{} % clear all footer fields
\renewcommand{\footrulewidth}{0.1pt}
\fancyfoot[L]{Truman Science Center}
\fancyfoot[C]{Stoichiometry Practice}
\fancyfoot[R]{5 November 2026}
%%
%
%% Title Set up
\styleTitle{Stoichiometry Practice}
%%
%
%%
%%% START TO DOCUMENT TEXT
%%
%
\textbf{For each of the following reactions...}
\begin{enumerate}[label=\alph*.]
    \item Write the balanced chemical equation. Use the table provided to keep track of how many of each group are on the product and reactant sides.
    \item Determine the Limiting Reagent.
    \item Calculate the Theoretical Yield in grams of the underlined product.
    \item Calculate the Percent Yield of the underlined product. 
\end{enumerate} 
\section{Abbreviations and Variable Names Used}

This section will be referring to the mock balanced chemical reaction below where lowercase letters are coefficients and uppercase letter are the compound.
\begin{gather*}
    \ce{aA + bB -> cC + dD}
\end{gather*}

\begin{table}[H]
\centering
\begin{tabular}{b{0.1\linewidth}|p{0.7\linewidth}} 
    $m_{atom}$ & Mass of one atom of this element \\
   \hline
    Count & The number of atoms of that element present in this compound \\
    \hline
    $\Sigma \ m_{atoms}$ & The collective mass of all atoms of that element in the molecule. Mathematically $\Sigma \ m_{atoms} = m_{atom} \cdot
    \text{Count}$ \\
    \hline
    MM & Molar Mass of the Molecule in $\qty{}{\frac{\gram}{\mole}}$ \\
    \hline
    $m_A$ & Given mass of chemical A \\
    \hline
    $MM_A$ & Molar Mass of chemical A\\
    \hline  
\end{tabular}
\end{table}
\section{The Strategy for Each Part}

\subsubsection{Part a}
You can only change the coefficients in front of each compound. The subscripts have to stay the same. Changing the subscripts changes what reaction is occurring. We only want to change how many of each molecule are being used at one time. 

\subsubsection{Part b} 
Convert the mass in grams of each reactant into a mass in grams of the underlined product. To convert between a measurement of one compound to another compound, you need to use the balanced chemical reaction between the two compounds.

To convert a mass in grams to an amount in moles, we need to calculate the Molar Mass of the compound. Using the elemental masses on the periodic table, add together the masses of all atoms in the compound. The molar mass of Carbon is $\qty{12.01}{\gram}$. The molar mass of Hydrogen is $\qty{1.01}{\gram}$. The molar mass of $\ce{CH4}$ which has 1 Carbon atom and 4 Hydrogen atoms is... 
\begin{gather*}
    \qty{12.01}{\gram} \ + \ (4 \ \cdot \ \qty{1.01 }{\gram}) \\ 
    \qty{16.05}{\gram}
\end{gather*}

This tells us how much of the product would be created if all of that reactant was used. Usually, one of the reactants will be able to make more product. The reactant that makes \textbf{less product} is the Limiting reactant. The general set up for this calculation in below. 
\begin{gather*} 
    m_A \ \cdot \frac{1}{MM_A} \cdot \frac{c \ \text{mol C}}{a \ \text{mol A}} \cdot MM_C = \boxed{m_C}
\end{gather*}

\subsubsection{Part c} 
The amount of product made by the limiting reactant calculated in part b is our theoretical yield. 

\subsubsection{Part d} 
Percent Yield compares the actual amount of product created (given in the question) with the theoretical maximum amount that \textbf{could} have been made. 
\begin{gather*}
    \text{Percent Yield} = 100\ \% \ \cdot \frac{\text{Actual Yield}}{\text{Theoretical Yield}}
\end{gather*}

\newpage

\textbf{1) $\qty{72.9}{\gram}$ of \ce{HCl} is mixed with \qty{148.2}{\gram} of \ce{Ca(OH)2} in an aqueous solution. This Acid-Base reaction yields \qty{103}{\gram} of \underline{\ce{CaCl2}}.}
\begin{table}[H]
\centering
\begin{tabular}{|c||c|c||c|c|} 
 % The multicolumn first row of the table is mostly for stylistic reasons
 % It starts with \phantom{Group xxxx} as a spacer so the reaction arrow is in the same place for each reaction
 \multicolumn{5}{c}{$\phantom{Group 1234567}
    \ce{HCl} \ + \ 
    \phantom{123} \ce{Ca(OH)2} \ \ce{->} \ 
    \phantom{123} \ce{H2O} \ + \ 
    \phantom{123} \ce{CaCl2}$} \\
 \hline
    & \multicolumn{2}{c||}{Reactants} & \multicolumn{2}{c|}{Products} \\
 \hline
 Group & Unbalanced & Balanced & Unbalanced & Balanced\\ 
 \hline\hline
   & & & & \\
 \hline
   & & & & \\
 \hline
  & & & & \\
 \hline
  & & & & \\[1ex] 
 \hline
\end{tabular}
\end{table}

\parskip =53mm

\textbf{2) $\qty{64.0}{\gram}$ of Carbon Dioxide and \qty{34}{\gram} of water forms glucose (\ce{C6H12O6}) and Oxygen gas. This reaction yields \qty{43.2}{\gram} of \underline{glucose}.} 
\begin{table}[H]
\centering
\begin{tabular}{|c||c|c||c|c|} 
 \multicolumn{5}{c}{$\phantom{Group 12345678901}
    \ce{CO2} \ + \ 
    \phantom{123} \ce{H2O} \ \ce{->} \ 
    \phantom{123} \ce{C6H12O6} \ + \ 
    \phantom{123} \ce{O2}$} \\
 \hline
    & \multicolumn{2}{c||}{Reactants} & \multicolumn{2}{c|}{Products} \\
 \hline
 Group & Unbalanced & Balanced & Unbalanced & Balanced\\ 
 \hline\hline
   & & & & \\
 \hline
   & & & & \\
 \hline
  & & & & \\
 \hline
  & & & & \\[1ex] 
 \hline
\end{tabular}
\end{table}


\textbf{3) $\qty{47.6}{\gram}$ of Magnesium Chloride is mixed with \qty{99.3}{\gram} of Lead (II) Nitrate in an aqueous solution. This double displacement reaction yields \qty{80.1}{\gram} of \underline{Lead (II) Chloride}.}
\begin{table}[H]
\centering
\begin{tabular}{|c||c|c||c|c|} 
 \multicolumn{5}{c}{$\phantom{Group 12345678}
    \ce{MgCl2} \ + \ 
    \phantom{123} \ce{Pb(NO3)2} \ \ce{->} \ 
    \phantom{123} \ce{PbCl2} \ + \ 
    \phantom{123} \ce{Mg(NO)3}$} \\
 \hline
    & \multicolumn{2}{c||}{Reactants} & \multicolumn{2}{c|}{Products} \\
 \hline
 Group & Unbalanced & Balanced & Unbalanced & Balanced\\ 
 \hline\hline
   & & & & \\
 \hline
   & & & & \\
 \hline
  & & & & \\
 \hline
  & & & & \\[1ex] 
 \hline
\end{tabular}
\end{table}

\textbf{4) The combustion of $\qty{47.6}{\gram}$ of \ce{C5H8O} in \qty{384.0}{\gram} of \ce{O2} yields \qty{398.7}{\gram} of gaseous \underline{\ce{CO2}} and water vapor.}
\begin{table}[H]
\centering
\begin{tabular}{|c||c|c||c|c|} 
 \multicolumn{5}{c}{$\phantom{Group 12345678}
    \ce{C5H8O} \ + \ 
    \phantom{123} \ce{O2} \ \ce{->} \ 
    \phantom{123} \ce{H2O} \ + \ 
    \phantom{123} \ce{CO2}$} \\
 \hline
    & \multicolumn{2}{c||}{Reactants} & \multicolumn{2}{c|}{Products} \\
 \hline
 Group & Unbalanced & Balanced & Unbalanced & Balanced\\ 
 \hline\hline
   & & & & \\
 \hline
   & & & & \\
 \hline
  & & & & \\
 \hline
  & & & & \\[1ex] 
 \hline
\end{tabular}
\end{table}



\textbf{5) $\qty{74}{\gram}$ of Carbon Dioxide is added to \qty{186.2}{\gram} of \ce{C9H10O2} as part of a synthesis reaction which yields \qty{239.2}{\gram} of \underline{\ce{C10H10O4}}.}
\begin{table}[H]
\centering
\begin{tabular}{|c||c|c||c|c|} 
 \multicolumn{5}{c}{$\phantom{Group 1234567890}
    \ce{CO2} \ + \ 
    \phantom{123} \ce{C9H10O2} \ \ce{->} \ 
    \phantom{123} \ce{C10H10O4} \ \phantom{+} \ 
    \phantom{123 \ce{CaCl2}}$} \\
 \hline
    & \multicolumn{2}{c||}{Reactants} & \multicolumn{2}{c|}{Products} \\
 \hline
 Group & Unbalanced & Balanced & Unbalanced & Balanced\\ 
 \hline\hline
   & & & & \\
 \hline
   & & & & \\
 \hline
  & & & & \\
 \hline
  & & & & \\[1ex] 
 \hline
\end{tabular}
\end{table}

\textbf{6) $\qty{147.8}{\gram}$ of Lithium Carbonate is mixed with \qty{636.8}{\gram} of Potassium Phosphate in an aqueous solution. This double displacement reaction yields \qty{152}{\gram} of \underline{Lithium Phosphate}.}
\begin{table}[H]
\centering
\begin{tabular}{|c||c|c||c|c|} 
 \multicolumn{5}{c}{$\phantom{Group 123456789}
    \ce{Li2CO3} \ + \ 
    \phantom{123} \ce{K3PO4} \ \ce{->} \ 
    \phantom{123} \ce{Li3PO4} \ + \ 
    \phantom{123} \ce{K2CO3}$} \\
 \hline
    & \multicolumn{2}{c||}{Reactants} & \multicolumn{2}{c|}{Products} \\
 \hline
 Group & Unbalanced & Balanced & Unbalanced & Balanced\\ 
 \hline\hline
   & & & & \\
 \hline
   & & & & \\
 \hline
  & & & & \\
 \hline
  & & & & \\[1ex] 
 \hline
\end{tabular}
\end{table}

\textbf{7) $\qty{1.51e24}{\molc}$ of Tetraphosphorous Hexoxide reacts with \qty{5.42e24}{\molc} of water. This reaction yields Phosphoric Acid and \qty{49.2}{\gram} of \underline{Phosphorous Trihydride}.}
\begin{table}[H]
\centering
\begin{tabular}{|c||c|c||c|c|} 
 \multicolumn{5}{c}{$\phantom{Group 1234567890}
    \ce{P4O6} \ + \ 
    \phantom{123} \ce{H2O} \ \ce{->} \ 
    \phantom{123} \ce{H3PO4} \ + \ 
    \phantom{123} \ce{PH3}$} \\
 \hline
    & \multicolumn{2}{c||}{Reactants} & \multicolumn{2}{c|}{Products} \\
 \hline
 Group & Unbalanced & Balanced & Unbalanced & Balanced\\ 
 \hline\hline
   & & & & \\
 \hline
   & & & & \\
 \hline
  & & & & \\
 \hline
  & & & & \\[1ex] 
 \hline
\end{tabular}
\end{table}


\textbf{8) $\qty{50.0}{\milli\liter}$ of $\qty{1.2}{\Molar}$ Aluminum Bromide is mixed with $\qty{30.0}{\milli\liter}$ of $\qty{1.33}{\Molar}$ Potassium Sulfate. This Double Displacement reaction yields \qty{4.5}{\gram} of \underline{Aluminum Sulfate}.}
\begin{table}[H]
\centering
\begin{tabular}{|c||c|c||c|c|} 
\multicolumn{5}{c}{$\phantom{Group 1234567890}
    \ce{AlBr3} \ + \ 
    \phantom{123} \ce{K2SO4} \ \ce{->} \ 
    \phantom{123} \ce{Al2(SO4)3} \ + \ 
    \phantom{123} \ce{KBr}$} \\
 \hline
    & \multicolumn{2}{c||}{Reactants} & \multicolumn{2}{c|}{Products} \\
 \hline
 Group & Unbalanced & Balanced & Unbalanced & Balanced\\ 
 \hline\hline
   & & & & \\
 \hline
   & & & & \\
 \hline
  & & & & \\
 \hline
  & & & & \\[1ex] 
 \hline
\end{tabular}
\end{table}

\textbf{9) $\qty{60.0}{\milli\liter}$ of $\qty{1}{\Molar}$ Phosphoric Acid is mixed with $\qty{100.0}{\milli\liter}$ of $\qty{0.800}{\Molar}$ Sodium Carbonate to create Sodium Phosphate, water, and gaseous carbon dioxide; it yields \qty{2.201}{\gram} of \underline{Carbon Dioxide}.}
\begin{table}[H]
\centering
\begin{tabular}{|c||c|c||c|c|} 
\multicolumn{5}{c}{$\phantom{Group 123456789012}
    \ce{H3PO4} \ + \ 
    \phantom{123} \ce{Na2CO3} \ \ce{->} \ 
    \phantom{123} \ce{Na3PO4} \ + \ 
    \phantom{123} \ce{CO2} \ + \
    \phantom{123} \ce{H2O}$}\\
 \hline
    & \multicolumn{2}{c||}{Reactants} & \multicolumn{2}{c|}{Products} \\
 \hline
 Group & Unbalanced & Balanced & Unbalanced & Balanced\\ 
 \hline\hline
   & & & & \\
 \hline
   & & & & \\
 \hline
  & & & & \\
 \hline
  & & & & \\[1ex] 
 \hline
\end{tabular}
\end{table}

\end{document}